\documentclass[11pt,letterpaper]{article}

\usepackage[margin=1in]{geometry}
\usepackage{amsmath,amssymb}
\usepackage{booktabs}
\usepackage{listings}
\usepackage{hyperref}
\usepackage{xcolor}
\usepackage[T1]{fontenc}
\usepackage{libertinus-type1}
\usepackage{microtype}

\lstset{
  basicstyle=\ttfamily\small,
  breaklines=true,
  frame=single,
  columns=fullflexible,
  keepspaces=true,
}

\hypersetup{
  colorlinks=true,
  linkcolor=blue!60!black,
  citecolor=blue!60!black,
  urlcolor=blue!60!black,
}

\title{Episodic Memory Without Learning:\\Ternary Neural Computation on Peripheral Hardware at 30 Microamps}
\author{Tripp Josserand-Austin\\EntroMorphic Research\\\texttt{tripp@entromorphic.com}}
\date{April 12, 2026}

\begin{document}
\maketitle

\begin{abstract}
We present a three-layer ternary neural computing system on a \$0.50 microcontroller that classifies wireless signals in peripheral hardware at 430\,Hz---96--100\% label-free accuracy on windowed classification (31--32/32 across 15 runs), 90.5--96.4\% per admitted packet---and accumulates a temporal model of what it has perceived---pattern-specific internal representations that develop from episodic memory retrieval over 120 seconds of live operation. The system draws approximately 30 microamps in autonomous mode. No floating point. No multiplication. No backpropagation. No training. Signatures enrolled from a 30-second observation window (sign of mean). CfC weights random and never updated.

The temporal model emerges from a 64-node Navigable Small World graph (2\,KB of LP SRAM) that stores and retrieves episodic snapshots of the system's perceptual state. The retrieval-and-blend mechanism alone---without any gate bias or weight learning---produces 8.5--9.7/80 MTFP divergence across 4 wireless signal patterns. The episodic memory is causally necessary: ablation (CfC without VDB blend) collapses pattern pairs that VDB feedback separates.

We also report two honest negative results from mechanisms designed to improve on the VDB baseline: (1)~agreement-weighted gate bias (kinetic attention) harms the representation in 2 of 3 runs (mean $-5.5 \pm 5.3$/80, $n=3$), and (2)~Hebbian weight learning produces no improvement ($+0.1 \pm 1.1$/80 MTFP, $n=3$). At $n=3$ neither mechanism is statistically distinguishable from zero; neither improves on the VDB baseline, and we report both under the same standard.

The system's power is in its episodic memory, not in learned weights or attentional bias. Classification is structurally separated from the temporal model by a zero-weight wall ($W_f[\text{hidden}] = 0$) that holds across all experiments. All results are from bare silicon.
\end{abstract}

\section{Introduction}

The dominant paradigm for neural computation on microcontrollers treats the device as a constrained deployment target---a model trained elsewhere, quantized, and loaded onto the chip to perform inference. The computation is designed in floating point and adapted to the hardware. The hardware is an inconvenience to be managed.

This paper describes a system designed the other way. The computation emerged from the hardware's native capabilities---ternary arithmetic on peripheral routing fabric---and the architecture grew from asking what structure those capabilities naturally support. The result is not a compressed version of a larger model. It is a system that could not exist in any other substrate, because its computational primitives (DMA descriptor chains, pulse counters, parallel I/O loopback) have no floating-point analog.

The system has three layers:

\textbf{Layer 1 (GIE):} A circular DMA chain streams ternary-encoded weight-input products through a parallel I/O peripheral configured in loopback mode. Two pulse counter units accumulate agreement and disagreement edges. An ISR decodes per-neuron dot products from cumulative differencing, applies a ternary CfC blend (UPDATE/HOLD/INVERT), re-encodes the updated hidden state into the DMA descriptors, and re-arms the parallel I/O byte counter. After initialization, the CPU computes zero dot products.

The chain runs at 430\,Hz with the CfC blend active. This is a direct measurement: the engine is started (which zeroes the loop counter), left to free-run for exactly one second, and the counter read---430 loops. It is the only clean rate measurement in our dataset, and every rate figure in this paper refers to it.

Classification runs in a different configuration, in which the blend is disabled (\texttt{gate\_threshold} $=$ \texttt{INT32\_MAX}, so no neuron fires and the hidden state is frozen) and the ISR skips the $\sim$20\,\textmu s re-premultiply-and-re-encode block. The loop rate in that configuration is necessarily higher, but we do not report a figure for it: the diagnostic our firmware printed was computed from counters spanning different windows and is not a valid rate (see Limitations). The classification results below are therefore reported without a rate attached.

\textbf{Layer 2 (LP core):} A 16\,MHz RISC-V ultra-low-power core, drawing approximately 30 microamps, runs a second ternary CfC network in hand-written assembly. It reads the GIE's hidden state, computes 32 ternary intersections against its own weight matrices, and maintains a 16-trit hidden state. A 64-node Navigable Small World graph provides episodic memory---storing 48-trit snapshots of the system's perceptual state and retrieving the nearest match for feedback blending into the LP hidden state.

\textbf{Layer 3 (HP core):} The 160\,MHz application core handles initialization, weight loading, ESP-NOW wireless reception, classification arbitration, and test orchestration. It is awake only when needed.

The operations used throughout are AND, popcount (byte lookup table), add, subtract, negate, branch, and shift. The operations absent are multiply, divide, floating point, gradients, and backpropagation.

\subsection{The Claims}

This paper demonstrates three things and honestly reports two failures:

\noindent\textbf{Demonstrated:}
\begin{enumerate}
\item Peripheral-hardware ternary dot products classify 4 wireless signal patterns label-free at 96--100\% on windowed classification (31--32/32 across 15 runs) and 90.5--96.4\% per admitted packet (TEST~11, TEST~14).
\item A 64-node episodic memory layer produces pattern-discriminative temporal states at 8.5--9.7/80 MTFP divergence from VDB feedback alone (TEST~12).
\item The episodic memory is causally necessary---ablation collapses what VDB feedback separates (TEST~13).
\end{enumerate}

\noindent\textbf{Failed:}
\begin{enumerate}
\setcounter{enumi}{3}
\item Agreement-weighted gate bias (kinetic attention) degrades LP representation at MTFP resolution in 2 of 3 runs: $-5.5 \pm 5.3$/80, $n=3$ (TEST~14).
\item Hebbian LP weight learning produces no improvement: $+0.1 \pm 1.1$/80, $n=3$ (TEST~15).
\end{enumerate}

Neither failure is statistically distinguishable from zero at $n=3$. What both establish is the same thing: neither mechanism improves on the VDB baseline, and the effort spent on them is better read as evidence for the baseline than against the mechanisms.

The VDB baseline IS the finding.

\section{Architecture}

\subsection{The Geometry Intersection Engine}

The GIE exploits a property of the ESP32-C6's peripheral fabric: the General DMA controller (GDMA) can stream data through the Parallel I/O transmitter (PARLIO) configured in 2-bit loopback mode, producing edge transitions on GPIO pins that are counted by two Pulse Counter (PCNT) units. By encoding ternary weight-input products as 2-bit pairs (positive = \texttt{01}, negative = \texttt{10}, zero = \texttt{00}), the PCNT accumulates agreement and disagreement edges across each neuron's 160-trit buffer. The dot product is $\text{agreement} - \text{disagreement}$.

The DMA descriptor chain is circular: 5 dummy neurons (all-zero buffers that absorb PCNT pipeline residue) followed by 64 real neurons (32 gate-pathway + 32 candidate-pathway), with the last descriptor pointing back to the first. Each neuron's buffer is followed by a 64-byte separator with EOF=1, which triggers a GDMA interrupt. The ISR captures cumulative PCNT values at each separator, decodes per-neuron dots via differencing, applies the CfC blend equation, re-encodes the hidden state into the DMA buffers, and re-arms PARLIO's byte counter.

The CfC blend equation in ternary:
\begin{align*}
f[n] &= \text{sign}(\text{dot}(W_f[n],\; [\text{input}\,|\,\text{hidden}])) \\
g[n] &= \text{sign}(\text{dot}(W_g[n],\; [\text{input}\,|\,\text{hidden}])) \\
h_{\text{new}} &= \begin{cases}
  h_{\text{old}} & f = 0 \quad \text{(HOLD)} \\
  g & f = +1 \quad \text{(UPDATE)} \\
  -g & f = -1 \quad \text{(INVERT)}
\end{cases}
\end{align*}

The three blend modes create non-gradient dynamics: HOLD provides inertia, UPDATE provides responsiveness, and INVERT creates oscillation and convergence resistance. Binary CfC has only UPDATE and HOLD. The third mode is a consequence of the ternary constraint.

\subsection{Classification}

\textbf{Input encoding (128 trits):} The ESP-NOW packet is encoded into 128 ternary trits:
\begin{itemize}
\item $[0..15]$ RSSI thermometer (16 trits)
\item $[16..23]$ Pattern ID one-hot (8 trits): \textbf{masked} under \texttt{MASK\_PATTERN\_ID\_INPUT=1}
\item $[24..87]$ Payload bits (64 trits): 8 bytes $\times$ 8 bits, bit$\to$trit
\item $[88..103]$ MTFP21 gap history (16 trits): inter-packet timing
\item $[104..127]$ Zeroed (24 trits)
\end{itemize}

\textbf{Label-free operation:} The \texttt{MASK\_PATTERN\_ID=1} flag zeros pattern\_id trits in the enrollment signatures. The \texttt{MASK\_PATTERN\_ID\_INPUT=1} flag zeros them in the runtime GIE input. Under both flags, no label information exists anywhere in the system. Classification relies entirely on payload content and inter-packet timing.

\textbf{TriX classification (ISR):} Four ternary signature vectors are installed as $W_f$ gate weights, with 8 neurons per pattern group. The ISR extracts per-group dot products and publishes the argmax. Accuracy: 31--32/32 (96--100\%) label-free on windowed classification; 90.5--96.4\% per admitted packet (Section~\ref{sec:classification}). The structural guarantee $W_f[\text{hidden}] = 0$ ensures classification is independent of the LP prior, gate bias, or temporal accumulation.

\textbf{Prior P1-P2 confusion (resolved):} With the original P2 payload, label-free accuracy was 71\%---P2 shared 48/64 payload trits with P1. The P2 payload was redesigned to produce a cross-dot of 29/96 (30\%), restoring label-free windowed accuracy to 31--32/32. This made the test fair, not easy.

\subsection{The LP Core: Geometric CfC + Episodic Memory}

The LP core runs a second CfC network with 16 hidden neurons, taking the GIE's 32-trit hidden state concatenated with its own 16-trit hidden state as a 48-trit input. Weights are random ternary, generated once at initialization and never updated.

The LP core's VDB stores up to 64 snapshots of the system's perceptual state---48-trit vectors comprising $[\text{gie\_hidden}\,|\,\text{lp\_hidden}]$. A Navigable Small World graph ($M=7$ neighbors) provides sub-linear approximate nearest-neighbor search. On each classification event (CMD~5), the LP core runs one CfC step, searches the VDB, and blends the best match's LP-hidden portion into the current LP hidden state.

The blend rule:
\begin{itemize}
\item Agreement ($h = \text{mem}$): no change
\item Gap fill ($h = 0$, $\text{mem} \neq 0$): adopt memory value
\item Conflict ($h \neq 0$, $\text{mem} \neq 0$, $h \neq \text{mem}$): set to 0 (HOLD)
\end{itemize}

The HOLD-on-conflict rule is the damper. A neuron in conflict reverts to the neutral state. This bounds the feedback loop---ternary values can change by at most 1 per step, and conflicting feedback produces inertia rather than oscillation.

\subsection{The MTFP Measurement Resolution}

The LP CfC computes 16 integer dot products per step (\texttt{lp\_dots\_f[16]}). The sign-quantized version (\texttt{lp\_hidden[16]} $= \text{sign}(\texttt{lp\_dots\_f})$) is the 16-trit vector used by the LP CfC's recurrence and the VDB. The MTFP-encoded version (5 trits per dot: sign + 2 exponent + 2 mantissa) is an 80-trit vector that preserves magnitude information.

The distinction is critical. Sign-space Hamming and MTFP Hamming can give \textbf{opposite results} for the same data (see Section~\ref{sec:sign-artifact}). All divergence measurements in this paper use MTFP unless otherwise noted.

\subsection{Memory Budget: 16\,KB}

The LP core's 16\,KB SRAM contains:
\begin{itemize}
\item Vector table: 128 bytes
\item Code (\texttt{.text}): $\sim$7,600 bytes
\item Popcount LUT (\texttt{.rodata}): 288 bytes
\item CfC state (\texttt{.bss}): $\sim$968 bytes
\item VDB nodes (\texttt{.bss}): 2,048 bytes (64 $\times$ 32 bytes)
\item VDB metadata: $\sim$80 bytes
\item Free for stack: $\sim$4,400 bytes (peak: 608 bytes during VDB search)
\end{itemize}

Every byte is specified. Every instruction is hand-written. No compiler decisions, no hidden register allocation, no spills.

\section{Results: What the Hardware Demonstrates}

\subsection{Classification: Windowed and Per-Packet (TEST 11, TEST 14)}
\label{sec:classification}

Four wireless signal patterns with distinct payloads and timing characteristics, classified from a 30-second enrollment window using sign-of-mean signatures. No labels in the input, no labels in the signatures, no labels anywhere in the system.

\textbf{Windowed classification (TEST~11): 31--32/32, i.e. 96--100\%.} Each sample is a majority vote over a window of $\sim$8--9 packets (276 packets $\to$ 32 windows). In the authoritative full-suite run both ISR TriX and CPU \texttt{core\_pred} are correct on all 32 windows. Across all 15 label-free runs captured after the P2-payload fix, the result is 32/32 in ten runs and 31/32 in five: the windowed figure is 96--100\%, not a uniform 100\%.

We report this figure with its window structure because the vote matters: on the same 32 windows, a trivial packet-rate baseline scores 28/32 = 87\%, so the windowed test has limited discriminating power.

\textbf{Per-admitted-packet classification (TEST~14): 90.5--96.4\%.} This is the TriX ISR prediction scored against sender ground truth, over $\sim$550 confirmations per condition, in the same firmware build and the same boot as the windowed result. The denominator is packets admitted by the novelty gate (CPU signature score $\ge$ 60), not all received packets:

\begin{table}[h]
\centering
\begin{tabular}{l rrrr r}
\toprule
Condition & P0 & P1 & P2 & P3 & Macro \\
\midrule
14A (no bias)  & 99.7 & 95.8 & 90.9 & \textbf{63.0} & \textbf{87.4} \\
14C-iso        & 99.0 & 95.1 & 92.3 & \textbf{51.8} & \textbf{84.6} \\
\bottomrule
\end{tabular}
\caption{TriX ISR per-class recall (\%) and macro average, label-free, per admitted packet, authoritative full-suite run. Overall accuracy is 90.5--96.4\%, but the sample distribution is heavily imbalanced (P0 $\approx$ 300 samples, P3 $\approx$ 46), so overall accuracy is dominated by P0. P3 is the weak class.}
\label{tab:per-packet}
\end{table}

The honest summary is therefore: the classifier is at or near perfect on windowed classification and imperfect per admitted packet, with essentially all of the error concentrated in P3. Section~\ref{sec:g3} shows that P3's neuron group is also the group that never fires the CfC gate---the two observations are likely the same phenomenon seen from opposite ends.

This is not a compressed model running inference. It is an enrollment-based classifier where the enrollment happens on the same hardware that runs the classification, using the same ternary operations, in 30 seconds of live observation.

\subsection{Temporal Context: 8.5--9.7/80 MTFP Divergence (TEST 12)}

With VDB feedback (CMD~5), 120 seconds of live wireless input from 4 cycling patterns produces the following LP MTFP divergence matrix:

\begin{table}[h]
\centering
\begin{tabular}{l rrrr}
\toprule
& P0 & P1 & P2 & P3 \\
\midrule
P0 & 0 & 5 & 6 & 8 \\
P1 & 5 & 0 & 9 & 13 \\
P2 & 6 & 9 & 0 & 10 \\
P3 & 8 & 13 & 10 & 0 \\
\bottomrule
\end{tabular}
\caption{LP MTFP divergence matrix (/80). Mean: 8.5/80 (full suite) to $9.7 \pm 0.6$/80 (3-rep measurement). All 6 pairs separated. P3 (ramp: incrementing payload at 10\,Hz) is the most distinctive.}
\label{tab:mtfp-matrix}
\end{table}

This is the LP CfC's 16 dot products, MTFP-encoded, accumulated over 120 seconds of episodic memory retrieval. The VDB stores snapshots of $[\text{gie\_hidden}\,|\,\text{lp\_hidden}]$ and blends the nearest match back into the LP state. Over time, the LP state develops pattern-specific magnitude profiles---different patterns produce different dot product distributions, visible in the MTFP encoding.

No labels drive this. No training signal. No gradient. The temporal model emerges from the interaction between the CfC projection (random but consistent) and the VDB episodic retrieval (stores what actually happened, blends it back in).

\subsection{Causal Necessity: VDB Feedback Required (TEST 13)}

CMD~4 (CfC + VDB search, no feedback blend) vs CMD~5 (CfC + VDB + feedback blend):
\begin{itemize}
\item CMD~4 P1-P2 sign Hamming: 1
\item CMD~5 P1-P2 sign Hamming: 2, MTFP Hamming: 9
\end{itemize}

VDB feedback contribution: +1 sign trit, +8 MTFP trits for the P1-P2 pair. The VDB is causally necessary.

\section{Honest Negatives: What We Tried and What We Learned}

\subsection{Kinetic Attention: Harmful at MTFP Resolution}

\textbf{Mechanism:} Agreement-weighted gate bias. When the LP prior agrees with the TriX prediction, lower the gate threshold for the predicted pattern's neuron group.

\textbf{Hypothesis:} Gate bias should amplify LP divergence by making the GIE more pattern-specific.

\begin{table}[h]
\centering
\begin{tabular}{l rrr}
\toprule
Run & 14A (no bias) & 14C (full bias) & Improvement \\
\midrule
1 & 9.8 & 10.2 & $+0.4$ \\
2 & 15.5 & 8.5 & $\mathbf{-7.0}$ \\
3 & 15.5 & 5.7 & $\mathbf{-9.8}$ \\
\midrule
Mean & 13.6 & 8.1 & $\mathbf{-5.5}$ \\
\bottomrule
\end{tabular}
\caption{MTFP divergence (/80) under kinetic attention. 3 runs, label-free, Test~14.}
\label{tab:kinetic}
\end{table}

\textbf{Statistical reading.} The effect is negative in 2 of 3 runs and positive in the third. Treating the three runs as independent replicates: mean $-5.47$, sd $5.27$, $t(2) = -1.80$, two-tailed $p \approx 0.21$. At $n=3$ the effect is \textbf{not statistically distinguishable from zero}. We apply the same standard here that we apply to Hebbian learning in Section~\ref{sec:hebbian}: neither mechanism separates from zero, and we decline to report one as a finding and the other as noise.

What the data does support, without qualification, is the weaker and more useful claim: \emph{kinetic attention does not improve on the VDB baseline}, and the largest single-run effects are degradations rather than improvements.

\textbf{Diagnosis:} The bias lowers the gate threshold for one neuron group. More neurons fire. A saturated GIE hidden state is more uniform across patterns, making the LP input less discriminative. The bias improves GIE sensitivity but degrades GIE discriminability.

\subsection{The Bias Reaches Only Three of Four Groups}
\label{sec:g3}

Per-group gate fire counts over 120\,s, identical in structure across all three runs:

\begin{table}[h]
\centering
\begin{tabular}{l rrrr}
\toprule
Condition & G0 & G1 & G2 & G3 \\
\midrule
14A (no bias)             & 82{,}608 & 27{,}825 & 36{,}320 & \textbf{0} \\
14C (full bias)           & 101{,}432 & 74{,}313 & 83{,}848 & \textbf{0} \\
14C-iso (bias after 60\,s) & 97{,}680 & 68{,}393 & 66{,}224 & \textbf{0} \\
\bottomrule
\end{tabular}
\caption{Per-group gate fires, run 1 (label-free). Group 3 never fires in any condition of any run.}
\label{tab:fires}
\end{table}

Neuron group 3 never fires the gate at \texttt{gate\_threshold} $=90$, with or without bias. The mechanism therefore modulates three of four pattern groups; the fourth is structurally inert, and no amount of bias reaches it.

This is a material limitation on the negative result. Kinetic attention was never tested on a full network. It is also the same group whose pattern (P3) carries essentially all of the per-packet classification error (Table~\ref{tab:per-packet}), which suggests P3's signature simply produces smaller dot magnitudes than the other three---a property of the enrolled signature, not of the bias mechanism.

The mechanism does fire where it can reach: per-group fire rate shifts exceed 10\% on G0--G2 in every run. The effect on the representation is negative on average and not distinguishable from zero. Reported honestly, with its scope.

\subsection{Hebbian Learning: Noise at 16 Neurons}
\label{sec:hebbian}

Three iterations, each addressing a limitation of the previous:

\begin{enumerate}
\item \textbf{VDB mismatch, f-only:} $+2.5$ with label in input, $-1.7$ genuinely label-free. The VDB error signal was exploiting the pattern\_id leaked through the GIE hidden state.

\item \textbf{TriX accumulator, f-only:} $-1.0$ label-free. Clean target but $\sim$50\% of errors were in the g-pathway, making f-pathway flips counterproductive.

\item \textbf{TriX accumulator, diagnosed f+g:} $+0.1 \pm 1.1$/80 MTFP at $n=3$. The diagnosis fixed the direction (harmful $\to$ neutral) but the mechanism produces no improvement.
\end{enumerate}

\textbf{Diagnosis:} The 16-neuron LP with random ternary weights produces 16 integer dot products. Each dot is the sum of $\sim$29 non-zero ternary products. Flipping one weight changes one dot by $\pm 2$---a tiny perturbation in a flat landscape with vast equivalence classes.

The VDB feedback blend, by contrast, directly injects episodic content into LP hidden. It doesn't need to find good weights---it bypasses the weights entirely.

\subsection{The Sign-Space Metric Can Mislead}
\label{sec:sign-artifact}

The kinetic attention finding was initially reported as positive ($+1.3$/16 in sign-space) before the MTFP numbers were examined.

\textbf{Run 2, the most dramatic case:}
\begin{itemize}
\item Sign-space 14A: 0.0/16 (all patterns $\to$ identical sign vector)
\item Sign-space 14C: 4.2/16 (bias ``broke the symmetry'')
\item MTFP 14A: 15.5/80 (patterns well-separated in magnitude)
\item MTFP 14C: 8.5/80 (bias crushed the magnitude diversity)
\end{itemize}

The LP states had identical signs but different magnitudes. Sign-space saw them as identical (Hamming~0). MTFP saw them as well-separated (15.5/80). The bias shuffled the signs (creating sign-space divergence) while crushing the magnitudes (destroying MTFP divergence). A net information loss that sign-space reported as a gain.

\textbf{The lesson:} When sign-space and MTFP disagree, trust MTFP.

\section{The Structural Wall}

The classification guarantee $W_f[\text{hidden}] = 0$ was verified across every experiment in the April 9--12 session, including label-free operation, Hebbian weight learning (LP $W_f$ and $W_g$ modified, GIE $W_f$ untouched), and multiple seeds.

TriX windowed accuracy is 31--32/32 across every label-free configuration we tested, and neither it nor per-admitted-packet accuracy (90.5--96.4\%) varies with the state of the temporal model. The structural wall---the guarantee that the classifier cannot be influenced by the temporal model---is the project's most robust finding.

The wall is verifiable by inspection, not merely by outcome. Every write to the hidden portion of $W_f$ in the firmware sets it to zero; the signature-install path and the online re-sign path both write only the input portion; and the LP Hebbian update touches the LP weight matrices exclusively. The one routine that could write $W_f$ at an arbitrary index is never called, and would not select a zero weight if it were.

This guarantee enables the honest negative results. We could test kinetic attention and Hebbian learning because the structural wall guaranteed that no matter how badly those mechanisms performed, they could not degrade classification. The classifier is immune to the prior. Note the scope precisely: the wall protects the \emph{classification} pathway ($W_f[\text{hidden}] = 0$), not the hidden state itself. Gate bias does modulate the hidden state, and therefore does reach the LP input---which is exactly the pathway on which kinetic attention was measured, and exactly why it could degrade the temporal representation while leaving classification untouched.

\section{Related Work}

\subsection{Ternary Neural Networks}

Ternary weight networks \cite{li2016ternary,zhu2016trained} quantize trained floating-point models to $\{-1, 0, +1\}$ for efficient inference. XNOR-Net \cite{rastegari2016xnor} and subsequent binary/ternary quantization methods \cite{zhang2018lqnets} treat ternary arithmetic as a compression technique---the model is designed in float, trained in float, and projected into low-precision weights for deployment. The computation is an approximation of the trained model.

The Reflex inverts this relationship. Its weights are generated ternary and never trained. There is no floating-point model being approximated. The ternary operations (AND, popcount, branch) are the native computation, not a quantized proxy for multiply-accumulate. The architectural consequences are different: ternary CfC has three blend modes (UPDATE/HOLD/INVERT) where binary has two, producing non-gradient dynamics that would not arise from quantizing a trained continuous model.

\subsection{Complementary Learning Systems}

McClelland, McNaughton, and O'Reilly \cite{mcclelland1995} proposed that intelligent systems require complementary fast (hippocampal) and slow (neocortical) learning systems. Kumaran et al.\ \cite{kumaran2016} updated the theory to emphasize replay-based consolidation. O'Reilly et al.\ \cite{oreilly2014} formalized the computational trade-offs between pattern separation (hippocampus) and pattern completion (neocortex).

The Reflex implements CLS with fixed neocortical weights---the CfC projection is random and never updated. This is closer to Benna and Fusi's \cite{benna2016} analysis of memory capacity in synaptic systems with constrained plasticity, where fixed random projections support episodic retrieval when the memory store has sufficient capacity. In our system, 64 NSW nodes in 2\,KB suffice for 4 patterns; the capacity bound is an open empirical question (see Section~\ref{sec:limitations}).

The VDB feedback blend implements memory-based prediction that parallels Lengyel and Dayan's \cite{lengyel2007} account of hippocampal contributions to decision-making: stored episodes directly inform current state estimation, bypassing learned statistical summaries. Our finding that Hebbian consolidation produces no improvement ($+0.1 \pm 1.1$) is consistent with their prediction that episodic systems outperform statistical learners when episodes are few and representations are low-dimensional.

See companion paper (Stratum~2) for the full CLS analysis, including the transition experiment and the permanent-hippocampus reframe.

\subsection{Neuromorphic and On-Device Computing}

Hardware neural networks on dedicated neuromorphic chips---Intel Loihi \cite{davies2018loihi}, BrainChip Akida \cite{vanschaik2023}, and IBM TrueNorth \cite{merolla2014truenorth}---implement spiking neural networks with on-chip learning rules (typically STDP). These systems achieve remarkable energy efficiency but require custom silicon. The Reflex uses commodity peripheral fabric (DMA, pulse counters, parallel I/O) on a \$0.50 general-purpose microcontroller, repurposing peripherals designed for motor control and communication as a neural substrate.

In the on-device inference space, TinyML frameworks \cite{banbury2021mlperf,warden2019tinyml} deploy quantized models (typically INT8) on microcontrollers using optimized runtimes (TFLite Micro, CMSIS-NN). These systems achieve inference but not temporal modeling---they classify each input independently, without accumulating a representation of what has been perceived over time. The Reflex's VDB episodic memory layer adds this temporal dimension at $\sim$30\,$\mu$A, below the sleep-mode power of most TinyML deployments that duty-cycle a full inference engine.

Approximate nearest-neighbor search on microcontrollers has been explored for resource-constrained retrieval \cite{jegou2011product}, but typically in service of classification or recommendation. The NSW graph in our system serves a different purpose: it provides episodic memory that feeds back into the neural dynamics, creating a recurrent loop between perception and memory that operates without gradient-based learning.

\subsection{Prior-Signal Separation}

The structural wall ($W_f[\text{hidden}] = 0$) that separates classification from temporal modeling relates to recent work on hallucination mitigation in language models. Retrieval-augmented generation \cite{lewis2020rag} provides external evidence but lacks architectural guarantees that the evidence cannot be overridden by the model's prior. See companion paper (Stratum~3) for the full five-component analysis of structural hallucination resistance.

\section{Limitations}
\label{sec:limitations}

\begin{enumerate}
\item \textbf{Single seed for most measurements.} The full-suite label-free dataset is one seed (0xCAFE1234). Kinetic attention was measured across 3 runs. Hebbian was replicated at $n=3$.

\item \textbf{Four patterns, one sender.} Classification uses 4 patterns with distinct payloads from a single ESP-NOW sender.

\item \textbf{JTAG attached.} All runs use USB-JTAG for serial output and power. The ``peripheral-autonomous'' claim requires UART-only verification with direct current measurement. Pending.

\item \textbf{MTFP metric is nonlinear, and /80 overstates its resolution.} MTFP Hamming counts differing trits in a nonlinear encoding. The 8.5--9.7/80 baseline should be interpreted as a proxy for dot-product diversity, not a calibrated distance. The encoding is also not of uniform capacity across its range: the fifth trit (the low mantissa trit) is computed as $\texttt{pos} \bmod \texttt{sub\_range}$ and is therefore constant ($-1$) for $|dot| \le 3$ and never takes $+1$ for $|dot| \le 24$. Only $|dot| \ge 25$ exercises all three states. With a 48-trit concatenated input at $\sim$40\% sparsity ($\sim$29 non-zero terms), typical dot magnitudes fall in exactly the low-capacity region. Over the full reachable range $[-48, +48]$ the encoder emits 75 distinct codes out of $3^5 = 243$ possible. Cross-condition comparisons in MTFP space remain valid---the same encoder is applied to every condition---but the ``/80'' denominator should not be read as 80 independent dimensions.

\item \textbf{Condition order is not counterbalanced.} Test~14 runs conditions in the fixed order 14A $\to$ 14C $\to$ 14C-iso within every run, and all three runs use that same order. Any drift across the $\sim$6-minute sequence---thermal, RF environment, board state---is confounded with condition and does not average out across replicates. The 60-second mid-run snapshot controls for representational maturity \emph{within} a condition, not for order \emph{between} conditions. The unstable 14A baseline across runs (9.8, 15.5, 15.5) is consistent with a drift term of this kind.

\item \textbf{Both negative results are underpowered.} At $n=3$, neither kinetic attention ($-5.5 \pm 5.3$) nor Hebbian learning ($+0.1 \pm 1.1$) separates from zero. They support the claim ``no improvement over the VDB baseline'' and do not support any claim about effect magnitude or reliability.

\item \textbf{The classification-mode loop rate is unmeasured.} Only the CfC-blend-active rate (430\,Hz) is cleanly measured. Our firmware also printed a ``TriX classifications at $N$\,Hz'' figure, but its numerator accumulated from the moment TriX was enabled while its denominator started later, at the beginning of the test window; the two spanned different intervals, so the figure overstates the rate by an unknown factor and is not cited here. The diagnostic has been corrected to measure both counters over one window, but no run under the corrected firmware exists yet.

\item \textbf{Transition experiment: single seed label-free.} The P1$\to$P2 transition experiment (TEST~14C) has been validated label-free on Seed~A. Multi-seed label-free TEST~14C has not been performed.

\item \textbf{16 neurons may be a ceiling.} The Hebbian null result may be specific to 16 neurons with 48-trit input. Wider LP or MTFP-targeted learning might enable weight learning.
\end{enumerate}

\section{Conclusion}

A ternary peripheral-fabric classifier on a \$0.50 microcontroller discriminates 4 wireless signal patterns---96--100\% label-free on windowed classification, 90.5--96.4\% per admitted packet---and builds a temporal context layer with 8.5--9.7/80 MTFP divergence, from episodic memory retrieval alone, without gate bias, without weight learning, without labels, without training, without floating point.

The system draws approximately 30 microamps. The classification is structurally separated from the temporal model by a zero-weight wall that holds across every experiment. The episodic memory (64 nodes, 2\,KB, NSW graph, 100\,Hz search) is causally necessary for LP divergence.

We tried to improve on the VDB baseline with two mechanisms: agreement-weighted gate bias and Hebbian LP weight learning. Gate bias is harmful---it saturates the GIE hidden state, reducing LP magnitude diversity. Hebbian learning is noise---the 16-neuron LP with random ternary weights is too coarsely quantized for single-trit flips to find useful structure.

The system's power is in its episodic memory. The temporal model is in the memories, not in the weights. The hardware computes the dot products. The memories hold the context. The structural wall guarantees they cannot corrupt each other.

\appendix

\section{Hardware Constants}

\begin{table}[h]
\centering
\begin{tabular}{ll}
\toprule
Parameter & Value \\
\midrule
Chip & ESP32-C6FH4 (QFN32) rev v0.2 \\
GIE rate & 430\,Hz (CfC blend active; classification-mode rate unmeasured) \\
LP core & 16\,MHz RISC-V, $\sim$30\,$\mu$A \\
LP CfC & 16 neurons, 48-trit input, fixed random weights \\
VDB & 64 nodes, NSW graph ($M=7$), 48-trit vectors \\
LP MTFP & 80 trits (5 per neuron: sign + 2 exp + 2 mant) \\
TriX & 4 patterns, 8 neurons/pattern, 32/32 windowed \\
Structural wall & $W_f[n][\texttt{CFC\_INPUT\_DIM}:] = 0$ for all neurons \\
Build flags & \texttt{MASK\_PATTERN\_ID=1}, \texttt{MASK\_PATTERN\_ID\_INPUT=1} \\
\bottomrule
\end{tabular}
\caption{System parameters.}
\end{table}

\section{Operations Used}

Gate pathway: AND, popcount (byte LUT), add, subtract, branch, shift. Candidate pathway: same operations, different weight matrix. Absent: multiply, divide, floating point, gradients, backpropagation.

\begin{thebibliography}{99}
\bibitem{li2016ternary}
F.~Li, B.~Zhang, and B.~Liu, ``Ternary weight networks,'' \textit{arXiv:1605.04711}, 2016.

\bibitem{zhu2016trained}
C.~Zhu, S.~Han, H.~Mao, and W.~J.~Dally, ``Trained ternary quantization,'' \textit{arXiv:1612.01064}, 2016.

\bibitem{rastegari2016xnor}
M.~Rastegari, V.~Ordonez, J.~Redmon, and A.~Farhadi, ``XNOR-Net: ImageNet classification using binary convolutional neural networks,'' \textit{ECCV}, pp.~525--542, 2016.

\bibitem{zhang2018lqnets}
D.~Zhang, J.~Yang, D.~Ye, and G.~Hua, ``LQ-Nets: Learned quantization for highly accurate and compact deep neural networks,'' \textit{ECCV}, pp.~365--382, 2018.

\bibitem{mcclelland1995}
J.~L.~McClelland, B.~L.~McNaughton, and R.~C.~O'Reilly, ``Why there are complementary learning systems in the hippocampus and neocortex,'' \textit{Psychological Review}, vol.~102, no.~3, pp.~419--457, 1995.

\bibitem{kumaran2016}
D.~Kumaran, D.~Hassabis, and J.~L.~McClelland, ``What learning systems do intelligent agents need? Complementary learning systems theory updated,'' \textit{Trends in Cognitive Sciences}, vol.~20, no.~7, pp.~512--534, 2016.

\bibitem{oreilly2014}
R.~C.~O'Reilly, R.~Bhatt, S.~Sungkhasettee, and K.~Shimizu, ``How sequential interactive processing within frontostriatal loops supports a continuum of habitual to controlled processing,'' \textit{Frontiers in Psychology}, vol.~5, p.~1541, 2014.

\bibitem{benna2016}
M.~K.~Benna and S.~Fusi, ``Computational principles of synaptic memory consolidation,'' \textit{Nature Neuroscience}, vol.~19, no.~12, pp.~1697--1706, 2016.

\bibitem{lengyel2007}
M.~Lengyel and P.~Dayan, ``Hippocampal contributions to control: The third way,'' \textit{Advances in Neural Information Processing Systems}, vol.~20, 2007.

\bibitem{davies2018loihi}
M.~Davies et al., ``Loihi: A neuromorphic manycore processor with on-chip learning,'' \textit{IEEE Micro}, vol.~38, no.~1, pp.~82--99, 2018.

\bibitem{vanschaik2023}
A.~van Schaik et al., ``The Akida neuromorphic processor,'' \textit{IEEE International Symposium on Circuits and Systems (ISCAS)}, 2023.

\bibitem{merolla2014truenorth}
P.~A.~Merolla et al., ``A million spiking-neuron integrated circuit with a scalable communication network and interface,'' \textit{Science}, vol.~345, no.~6197, pp.~668--673, 2014.

\bibitem{banbury2021mlperf}
C.~Banbury et al., ``MLPerf Tiny benchmark,'' \textit{Advances in Neural Information Processing Systems}, vol.~34, 2021.

\bibitem{warden2019tinyml}
P.~Warden and D.~Situnayake, \textit{TinyML: Machine Learning with TensorFlow Lite on Arduino and Ultra-Low-Power Microcontrollers}, O'Reilly Media, 2019.

\bibitem{jegou2011product}
H.~J\'{e}gou, M.~Douze, and C.~Schmid, ``Product quantization for nearest neighbor search,'' \textit{IEEE Trans.\ Pattern Analysis and Machine Intelligence}, vol.~33, no.~1, pp.~117--128, 2011.

\bibitem{lewis2020rag}
P.~Lewis et al., ``Retrieval-augmented generation for knowledge-intensive NLP tasks,'' \textit{Advances in Neural Information Processing Systems}, vol.~33, 2020.
\end{thebibliography}

\end{document}
